\section{Related Work}
\label{sec:related}

\subsection{Transparency Logs and Certificate Transparency}
Certificate Transparency (CT)~\cite{rfc6962} introduced the concept of
cryptographically verifiable, append-only logs for X.509 certificates. CT
requires that every publicly-trusted TLS certificate be submitted to at least
one log before issuance; browsers enforce this at the policy layer. Trillian~\cite{trillian}
generalises CT to arbitrary data structures via a Merkle-tree-backed
persistence layer. Both systems demonstrate that public, auditable logs are
technically feasible at scale. However, both assume voluntary operator
participation: the decision to submit is a runtime policy choice, not a
type-level constraint. BTV-Transparency closes this gap by making submission a
compile-time precondition for delivery.

\subsection{Linear Types for Resource Correctness}
Girard's linear logic~\cite{girard1987} treats propositions as resources
consumed exactly once. Wadler~\cite{wadler1990} connected this discipline to
programming language semantics, showing that linear types eliminate a class of
resource-management errors at compile time. The Rust programming
language~\cite{rust-book} operationalises linear ownership for memory safety.
BTV~\cite{btv2026} extended this paradigm to \emph{legal correctness},
encoding LGPD and EU AI Act obligations as linear resources. This paper extends
the same discipline one layer further, to \emph{persistence correctness}.

\subsection{AI Accountability and Policy-as-Code}
Policy-as-Code frameworks such as Open Policy Agent~\cite{opa} govern
infrastructure artefacts but assume a world of declarative configurations
rather than transient decision acts with legal consequences. Runtime audit
logs~\cite{nist-ai-rmf} can be truncated under load or dropped by a
misconfigured sink; they remain procedural promises, not structural invariants.
BTV demonstrated that the gap between runtime promises and compile-time
guarantees is bridgeable for in-memory correctness. BTV-Transparency extends
that bridge to the storage layer.
